\documentclass[10pt,a4paper,ngerman]{scrartcl}
\usepackage[utf8]{inputenc}
\usepackage[ngerman]{babel}
\usepackage[T1]{fontenc}
\usepackage{amsmath}
\usepackage{amsfonts}
\usepackage{amssymb}
\usepackage{makeidx}
\usepackage{graphicx}
\usepackage{epstopdf}
\usepackage[onehalfspacing]{setspace}
\usepackage{array}
\usepackage{upgreek}
\usepackage{float}
\usepackage{csquotes}
\usepackage{chemformula}
\usepackage{chemgreek}
\usepackage{chemmacros}
\chemsetup{modules=all}
\usepackage{tikz}
\usepackage{siunitx}
\sisetup{locale = DE,per-mode = symbol}
\usepackage{esvect}
\usepackage{geometry}
\usepackage{pdfpages}
\usepackage[font=small]{caption}
%\usepackage[version=4]{mhchem}
\usepackage[backend=biber, style=chem-angew]{biblatex}
\addbibresource{Polyaddition_Polykond.bib}
\usepackage{tabularx, booktabs, multirow}
\usepackage[pdfborder={0 0 0}]{hyperref}
\usepackage{scrlayer-scrpage}%Header für die KOMA-script -Klasse
%\usepackage{hyperref}
\usepackage{pgfplots}
\usepackage{textgreek}


\newcommand{\tildenu}{{\fontencoding{LGR}\selectfont\accperispomeni\textnu}}
\captionsetup{format=plain}
\chemsetup[phases]{pos=sub}
\chemsetup[redox]{pos=top}
\chemsetup[redox]{align=center}
%\chemsetup[redox]{explicit-sign=true}
%\chemsetup[redox]{roman=false}
\newcommand{\celsius}{^{\circ}\mathrm{C}}
\parindent0pt
\sloppy
\DeclareChemPhase{\s}{s}
\DeclareChemPhase{\l}{l}
\DeclareChemPhase{\g}{g}
\renewcaptionname{ngerman}{\figurename}{Abb.}
\renewcaptionname{ngerman}{\tablename}{Tab.}
\geometry{bottom=100pt} \geometry{top=80pt}
\newcommand{\m}{\mathrm{m}}
\newcommand{\K}{\mathrm{K}}
\newcommand{\J}{\mathrm{J}}
\newcommand{\gr}{\mathrm{g}}
\newcommand{\kg}{\mathrm{kg}}
\newcommand{\mol}{\mathrm{mol}}
\newcommand{\Pa}{\mathrm{Pa}}
\newcommand{\ad}{\mathrm{ad}}
\newcommand{\de}{\mathrm{de}}
\newcommand{\gmol}{\mathrm{\frac{g}{mol}}}
\newcommand{\JKmol}{\mathrm{\frac{J}{K \cdot mol}}}
\newcommand{\kgmss}{\mathrm{\frac{kg}{m \cdot s^2}}}
\newcommand{\loge}{\mathrm{ln}}


\begin{document}
\begin{titlepage}
	\vspace*{2,5cm}
	\begin{center}
		\begin{LARGE}
			\textbf{Polymer Chemie}\\
			\textbf{Gelpermeationschromatographie (GPC)}\\
		\end{LARGE}
		\vspace{1cm}
		\textbf{Universität Stuttgart}\\
		\vspace*{1,5cm}
		\begin{tabular}{lp{7,5cm}l}
			Author 1:     &Laurens Viehoff\\
            Author 2:     &Felix Roth\\   
			E-Mail 1:      &st183688@stud.uni-stuttgart.de\\
            E-Mail 2:      &st188157@stud.uni-stuttgart.de  \\
			Student-ID 1: & 3676761   \\ 
            Student-ID 2: & 3711192 \\ \\
			Assistentin:  & Klaus Dirnberger  \\ \\ 
		\end{tabular}
	\end{center}

    
\end{titlepage}
\thispagestyle{empty}
\renewcommand{\thepage}{\arabic{page}}
\newpage
\tableofcontents
\setcounter{page}{1}
\renewcommand{\thepage}{\arabic{page}}
\newpage

\section{Theorie}

	\subsection{Grundprinzipien}
		
		Bei der Gelpermeationschromatographie (GPC) werden die Molmassen der Polymere Chromatographisch aufgetrennt und bestimmt. Dabei Durchläuft 
		das Polymer eine stationäre Phase, die aus einem porösen Gel besteht. Beim passieren der Polymere durch die stationäre Phase, gelangen 
		Polymerknäuel in die Poren der stationären Phase und werden somit zurückgehalten. Dabei können sehr große Polymerknäuel nicht in die Poren 
		gelangen, da sie zu groß sind. Diese werden nicht zurückgehalten und werden als erstes detektiert aber nicht aufgetrennt. Anschließend werden 
		die kleinen Polymerknäuel durch ein Konzentrationsgradienten aus den Poren gezogen und detektiert. Davon werden allerdings nur sehr kleine 
		Polymerknäuel beeinflusst, da sie ein größeres Permeationsvermögen besitzen und als zweites Detektiert. Diese beiden Grenzen bilden das Fenster 
		in denen die Polymerknäuel aufgetrennt werden können. So muss je nach Größe der Polymerknäuel eine andere Porengröße gewählt werden, damit dieses 
		Auftrennungsfenster optimal genutzt werden kann. Bei vorgegebener Porengröße werden die Polymere mit größen zwischen der oberen und unteren 
		Grenzen von groß nach klein aufgetrennt und die Intensität der Detektion ist dabei dann proportional zur Menge an Polymer dieser Molmasse. \\
		Die Größe der Polymerknäuel hängt dabei von der Molmasse und dem Solvatationsvermögen der Polymere im gewählten Lösungsmittel ab. Werden 
		Beispielsweise zwei Polymere mit identischer Molmasse in einem bestimmten Lösungsmittel gelöst, sind die Wechselwirkungen zwischen dem einem 
		Polymer zum Lösungsmittel stärker und lassen dieses weiter aufknäulen. Somit würde dieses auch eine andere Retentionszeit besitzen und früher 
		Detektiert werden. \\

	\subsection{Aufbau der GPC}

		


\section{Auswertung}



\section{Durchführung}



\section{Fragen}
\printbibliography

\end{document}